\section{Záver}

Cieľom tímového projektu bolo navrhnúť a realizovať inteligentný IoT systém 
na riadenie teploty kvapaliny pomocou Peltierovho článku. Výsledkom riešenia 
je laboratórny model, ktorý prepája hardvérovú časť, mikrokontrolér ESP32, 
lokálny Flask server, webové rozhranie, dátové úložisko a cloudovú platformu ThingsBoard.

Hardvérová časť systému bola navrhnutá ako uzavretý hydraulický okruh s vodnou pumpou, 
výmenníkom tepla, Peltierovým článkom, chladičom s ventilátorom a snímačom teploty. 
Pri návrhu bolo potrebné zohľadniť najmä prúdové zaťaženie výkonových prvkov, 
stabilitu napájania, bezpečné oddelenie elektroniky od vodného okruhu a dostatočné 
chladenie teplej strany Peltierovho článku. Súčasťou hardvérového riešenia bolo 
aj lokálne zobrazovanie údajov na LCD displeji a postupné upravovanie zapojenia 
podľa potrieb testovania.

Softvérová časť zabezpečuje komunikáciu medzi ESP32 a serverom, spracovanie aktuálneho stavu systému, 
realtime zobrazovanie údajov vo webovom rozhraní a ukladanie historických meraní. 
Na komunikáciu medzi webovým rozhraním a serverom bola použitá WebSocket komunikácia, 
zatiaľ čo dáta zo zariadenia sú prijímané cez REST API. Merania sa ukladajú do MySQL databázy 
a zároveň aj do záložného súboru vo formáte JSON, čo umožňuje spätné zobrazenie a 
vyhodnocovanie nameraných priebehov.

Dôležitou súčasťou projektu bola aj integrácia s platformou ThingsBoard. 
Do cloudového dashboardu sa odosielajú telemetrické údaje, ako aktuálna teplota, 
požadovaná teplota, regulačná odchýlka, PWM hodnota pumpy, PWM hodnota Peltierovho článku, 
stav systému a aktívny režim. Dashboard zároveň umožňuje základné vzdialené ovládanie systému, 
prepínanie režimov a manuálne nastavovanie PWM hodnôt. Okrem troch regulačných módov požadovaných 
zadaním bol doplnený aj štvrtý manuálny režim, ktorý umožňuje priame ovládanie pumpy aj Peltierovho článku.

Projekt ukazuje prepojenie viacerých oblastí mechatroniky, automatického riadenia a IoT technológií. 
Riešenie spája reálnu tepelnú sústavu, výkonovú elektroniku, meranie teploty, webovú aplikáciu, 
databázové ukladanie dát a cloudový monitoring do jedného funkčného celku. Navrhnutý systém vytvára 
vhodný základ na ďalšie testovanie regulačných algoritmov, identifikáciu tepelnej sústavy a overovanie 
správania systému pri rôznych režimoch riadenia.

